\documentclass[english,hebrew]{article}
\usepackage[HE8,T1]{fontenc}
\usepackage[utf8]{inputenc}
\setcounter{secnumdepth}{-2}
\setlength{\parindent}{0bp}
\usepackage{amsmath}
\usepackage{amsthm}
\usepackage{amssymb}

\makeatletter
\theoremstyle{plain}
\newtheorem{thm}{\protect\theoremname}
\theoremstyle{definition}
\newtheorem{problem}[thm]{\protect\problemname}
\theoremstyle{definition}
\newtheorem*{sol*}{\protect\solutionname}

\@ifundefined{date}{}{\date{}}
\usepackage{graphicx}
\usepackage{geometry}
\newcommand{\limi}{\underset{n\rightarrow\infty}{\lim}}
\newcommand{\limxz}{\underset{x\rightarrow x_{0}}{\lim}}
\newcommand{\limfxz}{\underset{x\rightarrow x_{0}}{\lim}f(x)}
\newcommand{\limfxm}{\underset{x\rightarrow x_{0}^{-}}{\lim}f(x)}
\newcommand{\limfxp}{\underset{x\rightarrow x_{0}^{+}}{\lim}f(x)}
\newcommand{\limxm}{\underset{x\rightarrow x_{0}^{-}}{\lim}}
\newcommand{\limxp}{\underset{x\rightarrow x_{0}^{+}}{\lim}}
\newcommand{\sqrm}{M_{m\times m}(\mathbb{F})}
\newcommand{\seqa}{(a_{n})_{n=1}^{\infty}}
\newcommand{\seqx}{(x_{n})_{n=1}^{\infty}}
\newcommand{\funcdr}{f:D\rightarrow\mathbb{R}}
\newcommand{\der}{\limxz\frac{f(x)-f(x_{0})}{x-x_{0}}}
\usepackage[all]{pdfprivacy}

\makeatother

\usepackage{babel}
\addto\captionsenglish{\renewcommand{\problemname}{Problem}}
\addto\captionsenglish{\renewcommand{\solutionname}{Solution}}
\addto\captionsenglish{\renewcommand{\theoremname}{Theorem}}
\addto\captionshebrew{\renewcommand{\problemname}{בעיה}}
\addto\captionshebrew{\renewcommand{\solutionname}{פתרון}}
\addto\captionshebrew{\renewcommand{\theoremname}{משפט}}
\providecommand{\problemname}{בעיה}
\providecommand{\solutionname}{פתרון}
\providecommand{\theoremname}{משפט}

\begin{document}
\newgeometry{left=1cm, right=1cm, top=1cm, bottom=1cm}
\title{סמסטר {\beginL 2027\endL}א, בוחן {\beginL \vspace{-1em}1\endL}}
\maketitle
\begin{abstract}
בוחן {\beginL 1\endL} בקורס אינפי {\beginL 2\endL}, סמסטר תשפ\textquotedblright ז
א'. השאלות נכתבו על ידי צוות הקורס באוניברסיטה העברית. הפתרון הוא
שלי ולא רשמי. \\
\end{abstract}
\begin{problem}
מצאו את כל הפונקציות הקדומות של הפונקציה \L{$f(x)=e^{2x}\cos\left(e^{x}\right)$}.
\end{problem}

\begin{sol*}
מבין הכלים שלמדנו, אין כאן גזירה ישירה, אין כאן מנת פולינומים ולינאריות
לא תעזור כאן. 

שתי השיטות הרלוונטיות הן אינטגרציה בהצבה או אינטגרציה בחלקים. 

בשביל אינטגרציה בחלקים צריך ביטוי שנרצה לגזור כפול ביטוי שנוח לנו
לטנגרל, וזה לא המצב.

מכאן, כנראה שהכלי הרלוונטי הוא אינטגרציה בהצבה.

בשביל אינטגרציה בהצבה צריך ביטוי שהנגזרת שלו מתחבאת בביטוי, וזה בדיוק
מה שקורה.

נבחר \L{$t=e^{x}$} )בבחנים אני כותב \L{$t$} כי \L{$u$} נראה כמו
{\beginL 4\endL}(. אז \L{%
\[
\frac{dt}{dx}=e^{x}\Rightarrow dt=e^{x}dx
\]
}%
\L{%
\[
\int e^{2x}\cos\left(e^{x}\right)dx=\int\left(e^{x}\right)^{2}\cos\left(e^{x}\right)dx=\int e^{x}\cos\left(e^{x}\right)e^{x}dx
\]
}%
\L{%
\[
=\int t\cos\left(t\right)dt
\]
}%

כאן יש ביטוי שנרצה לגזור כפול ביטוי שנוח לנו לטנגרל, לכן נשתמש באינטגרציה
בחלקים\L{%
\[
=\int t\cdot\left(\sin(t)\right)^{'}dt=t\sin(t)-\int1\cdot\sin(t)dt=t\sin(t)+\cos(t)+C
\]
}%
\L{%
\[
=e^{x}\sin\left(e^{x}\right)+\cos\left(e^{x}\right)+C
\]
}%
\end{sol*}
\begin{problem}
נסתכל על הפונקציה\L{%
\[
\tanh(x)=\frac{e^{x}-e^{-x}}{e^{x}+e^{-x}}
\]
}%

האם יש לה גבול כאשר \L{$x\rightarrow\infty$}? אם כן, מהו?

האם יש לה גבול כאשר \L{$x\rightarrow-\infty$}? אם כן, מהו?

האם הפונ' מונוטונית?
\end{problem}

\begin{sol*}
נשים לב כי היות ו-\L{$\forall x\in\mathbb{R}:e^{x}\in\left(0,\infty\right)$}
מתקיים גם\L{%
\[
\tanh(x)=\frac{e^{x}-\frac{1}{e^{x}}}{e^{x}+\frac{1}{e^{x}}}=\frac{e^{2x}-1}{e^{2x}+1}=\frac{1-\frac{1}{e^{2x}}}{1+\frac{1}{e^{2x}}}
\]
}%

ניזכר כי \L{$\underset{x\rightarrow\infty}{\lim}e^{x}=\infty$}, לכן
לפי אריתמטיקה של גבולות במובן הרחב \L{%
\[
\underset{x\rightarrow\infty}{\lim}\frac{1}{e^{2x}}=\underset{x\rightarrow\infty}{\lim}\frac{1}{e^{x}}\cdot\underset{x\rightarrow\infty}{\lim}\frac{1}{e^{x}}=0\cdot0=0
\]
}%
 ומכאן\L{%
\[
\underset{x\rightarrow\infty}{\lim}\tanh(x)=\underset{x\rightarrow\infty}{\lim}\frac{1-\frac{1}{e^{2x}}}{1+\frac{1}{e^{2x}}}=\frac{1-0}{1+0}=1
\]
}%

בנוסף ניזכר כי \L{$\underset{x\rightarrow-\infty}{\lim}e^{x}=0$}
לכן לפי אריתמטיקה\L{%
\[
\underset{x\rightarrow-\infty}{\lim}\left(\frac{e^{2x}-1}{e^{2x}+1}\right)=\left(\frac{0-1}{0+1}\right)=-1
\]
}%

למדנו כי הפונקציה \L{$e^{x}$} היא פונקציה רציפה.

לפי אריתמטיקה של פונ' רציפות, גם הפונקציות \L{$e^{2x}-1,e^{2x}+1$}
רציפות.

בנוסף, הפונ' \L{$e^{2x}+1$} אף פעם לא מתאפסת לפי מה שציינו על הטווח
של \L{$e^{x}$}, לכן מאריתמטיקה הפונ' כולה רציפה.

יהיו \L{$x_{1},x_{2}\in\mathbb{R}$} כך ש-\L{%
\[
\tanh(x_{1})=\tanh(x_{2})
\]
}%

אז\L{%
\[
\frac{e^{2x_{1}}-1}{e^{2x_{1}}+1}=\frac{e^{2x_{2}}-1}{e^{2x_{2}}+1}
\]
}%
\L{%
\[
\underset{e^{2t}>0}{\iff}\left(e^{2x_{1}}-1\right)\left(e^{2x_{2}}+1\right)=\left(e^{2x_{1}}+1\right)\left(e^{2x_{2}}-1\right)
\]
}%
\L{%
\[
\iff e^{2(x_{1}+x_{2})}+e^{2x_{1}}-e^{2x_{2}}-1=e^{2(x_{1}+x_{2})}-e^{2x_{1}}+e^{2x_{2}}-1
\]
}%
\L{%
\[
\iff e^{2x_{1}}-e^{2x_{2}}=-e^{2x_{1}}+e^{2x_{2}}
\]
}%
\L{%
\[
\iff2e^{2x_{1}}=2e^{2x_{2}}
\]
}%
ניזכר שהפונ' \L{$\exp$} היא מונוטונית ממש ורציפה, לכן היא חח\textquotedblright ע
ומכאן \L{$x_{1}=x_{2}$} כלומר גם הפונ' \L{$\tanh$} היא חח\textquotedblright ע.

מכך שפונ' רציפה וחח\textquotedblright ע היא מונוטונית ממש, הפונ' \L{$\tanh$}
מונטונית ממש וסיימנו.
\end{sol*}

\end{document}
